\section{Sample Path Ergodicity}
\label{sec:A.6}

Why do we care about ergodicity of Markov chains in a Monte Carlo simulation course?  Markov chain
Monte Carlo (MCMC) algorithms draw samples from a Markov chain which has as invariant a prescribed
target distribution $\pi = f$.  Ultimately, we are interested in understanding the \emph{sample
path ergodicity} of MCMC chains, that is, we want to determine if ergodic averages
\[
  \frac{1}{N} \sum_{n=1}^{N} h(X^{(n)})
\]
computed from MCMC samples accurately approximate expectations $\calI_f[h]$ with respect to the
target.  These samples are \emph{not} independent, and consequently the efficiency of the ergodic
average estimator will depend on the correlations between them.

First, we have the following law of large numbers.

\begin{theorem}[Law of Large Numbers]
\label{thm:A.32}
Let $h : E \to \R$ and let $\mathbb{X} = \{X_n\}_{n=0}^\infty$ be an ergodic Markov chain with
stationary distribution $\pi$.  Then, $\pi$-almost surely,
\[
  \frac{1}{N} \sum_{n=1}^{N} h(X_n) \xrightarrow{N \to \infty} \int_E h(x)\, \pi(x)\, dx
  \equiv \calI_\pi[h].
\]
\end{theorem}

We also have a central limit theorem, which requires geometric ergodicity of the chain.

\begin{definition}[Geometric and Uniform Ergodicity]
\label{def:A.33}
An ergodic Markov chain with invariant distribution $\pi$ is geometrically ergodic if there exist
a non-negative function $M$ with $\Expect_\pi[M(X)] < \infty$ and $0 < r < 1$ such that
\[
  d_{\mathrm{TV}}(P^n(x, \cdot), \pi) \leq M(x)\, r^n
\]
for all $x$ and all $n$.  If the function $M$ is bounded above, then the chain is called
uniformly ergodic.
\end{definition}

Note, as an example, that we have proved that any irreducible and aperiodic Markov chain on a
finite state space is uniformly ergodic.  The following result shows that sample path ergodicity
can be deduced from geometric ergodicity.  In Chapter~4, we discuss the uniform and geometric
ergodicity of several MCMC algorithms, thus guaranteeing in particular their sample path
ergodicity.

\begin{theorem}[Central Limit Theorem]
\label{thm:A.34}
With the same notation as in the previous theorem, suppose further that $\mathbb{X}$ is
geometrically ergodic and that, for some $\epsilon > 0$,
\[
  \Expect_{Y \sim \pi}[h(Y)^{2+\epsilon}] < \infty.
\]
Then,
\[
  \sqrt{N}\!\left(\frac{1}{N} \sum_{n=1}^{N} h(X_n) - \calI_\pi[h]\right)
  \cvd \Normal(0, \tau^2),
\]
where $\tau^2$ is the integrated autocorrelation time, defined below in Lemma~\ref{lem:A.36}.
\end{theorem}

The rest of this appendix is devoted to defining and understanding the integrated autocorrelated
time, which may be used to assess the relative efficiency of different MCMC algorithms.

Unsurprisingly, the asymptotic variance $\tau^2$ in the above central limit theorem arises as the
limiting variance of ergodic averages.  We define
\[
  \frac{\tau_N^2}{N} := \Var\!\left[\frac{1}{N} \sum_{n=1}^{N} h(X_n)\right]
\]
and note that $\tau_N^2$ (and similarly $\sigma^2$ and $\tau$ below) depends on $h$ but we omit
said dependence from our notation.  The value $\tau_N^2$ quantifies the amount of correlation
between the variables $X_n$'s.

\begin{definition}[Autocovariance and Autocorrelation]
\label{def:A.35}
At stationarity, the autocovariance of lag~$k$ of the time series $\{h(X_n)\}_{n=0}^\infty$ is
defined as
\[
  \gamma_k = \mathrm{Cov}(h(X_0), h(X_k)), \qquad k \geq 1.
\]
We also define $\sigma^2 = \mathrm{Cov}(h(X_0), h(X_0)) = \Var_{X_0 \sim \pi}[h(X_0)]$.  The
autocorrelation of lag~$k$ is
\[
  \rho_k = \frac{\gamma_k}{\sigma^2}, \qquad k \geq 0.
\]
\end{definition}

\begin{lemma}
\label{lem:A.36}
At stationarity,
\[
  \tau_N^2 = \sigma^2 \left[1 + 2\sum_{k=1}^{N-1} \frac{N-k}{N}\,\rho_k\right].
\]
As $N \to \infty$,
\[
  \tau_N^2 \to \tau^2 := \sigma^2 \left[1 + 2\sum_{k=1}^{\infty} \rho_k\right],
\]
where $\tau^2$ is called the integrated autocorrelation time.
\end{lemma}

\begin{proof}
A calculation shows that
\begin{align*}
  \frac{\tau_N^2}{N}
    &= \Var\!\left[\frac{1}{N}\sum_{n=1}^{N} h(X_n)\right] \\
    &= \frac{1}{N^2}\left[\sum_{n=1}^{N} \Var[h(X_n)]
       + 2\sum_{n=1}^{N-1}\sum_{k>n} \mathrm{Cov}(h(X_n), h(X_k))\right] \\
    &= \frac{1}{N^2}\left[N\sigma^2
       + 2\sum_{n=1}^{N-1}\sum_{k=1}^{N-n} \mathrm{Cov}(h(X_n), h(X_{n+k}))\right] \\
    &= \frac{\sigma^2}{N}\left[1
       + \frac{2}{N\sigma^2}\sum_{n=1}^{N-1}\sum_{k=1}^{N-n} \gamma_k\right] \\
    &= \frac{\sigma^2}{N}\left[1
       + \frac{2}{N}\sum_{n=1}^{N-1}\sum_{k=1}^{N-n} \rho_k\right] \\
    &= \frac{\sigma^2}{N}\left[1
       + \frac{2}{N}\sum_{k=1}^{N-1}\sum_{n=1}^{N-k} \rho_k\right] \\
    &= \frac{\sigma^2}{N}\left[1
       + 2\sum_{k=1}^{N-1} \frac{N-k}{N}\,\rho_k\right]. \qedhere
\end{align*}
\end{proof}

Note that if $X_n \iid \pi$, then
\[
  \Var\!\left[\frac{1}{N}\sum_{n=1}^{N} h(X_n)\right] = \frac{\sigma^2}{N}.
\]
Therefore, computing an ergodic average with positively autocorrelated samples leads to higher
variance than computing it with an i.i.d.\ sample.  An intuitive explanation is that positively
correlated random variables have redundant information so are less informative than i.i.d.\ random
variables.  On the other hand, if the correlations are negative, ergodic averages may be more
accurate than a direct Monte Carlo estimator with i.i.d.\ samples.
